\section{Algorithm and Context Specification}
\label{app:algorithm_specification}

\Cref{fig:allocator_policy_architecture,alg:joint_loop} distinguish inter-route budgeting from model-specific within-route allocation selection.

\begin{figure*}[t]
\centering
\resizebox{0.92\textwidth}{!}{%
\begin{tikzpicture}[>=Latex,node distance=0.55cm and 0.55cm,box/.style={rectangle,draw=black!70,rounded corners=2pt,align=center,minimum height=0.75cm,text width=2.4cm,fill=black!3},arr/.style={thick,->}]
\node[box] (budget) {Primary physical budget\\$Q=35$};
\node[box,right=of budget] (allocator) {Inter-route allocator};
\node[box,right=of allocator] (routes) {Route budgets\\$T_1,\ldots,T_R$};
\node[box,right=of routes] (sets) {Feasible sets\\$\mathcal{X}_r(T_r)$};
\node[box,right=of sets] (model) {Model-specific\\route selection};
\node[box,right=of model] (action) {Model-specific within-route\\allocation action};
\node[box,below=0.85cm of action] (reward) {Threat-conditioned\\reward};
\node[box,left=of reward] (update) {Model-specific\\update};
\node[box,left=of update] (stats) {Replay and\\route statistics};
\draw[arr] (budget)--(allocator); \draw[arr] (allocator)--(routes); \draw[arr] (routes)--(sets); \draw[arr] (sets)--(model); \draw[arr] (model)--(action); \draw[arr] (action)--(reward); \draw[arr] (reward)--(update); \draw[arr] (update)--(stats);
\draw[arr,dashed] (stats.west) -| node[pos=0.25,below,align=center]{if allocator update\\is enabled and due} (allocator.south);
\end{tikzpicture}%
}
\caption{Allocator--policy architecture. The 35-qubit budget is specific to the primary matched evaluation, not a universal framework limit. The dashed feedback path is conditional; no universal allocator-update cadence is asserted.}
\label{fig:allocator_policy_architecture}
\end{figure*}

\begin{algorithm}[t]
\caption{Shared routing and within-route allocation decision loop}
\label{alg:joint_loop}
\begin{algorithmic}[1]
\scriptsize
\Require topology $G$, budget $Q$, allocator $A$, policy $\pi$, threat process, replay setting
\State $\mathbf{T}\gets A.\Call{Initialize}{Q,\mathcal{P}}$
\ForAll{$P_r\in\mathcal{P}$}
  \State $\mathcal{X}_r(T_r)\gets\Call{FeasibleAllocations}{P_r,T_r}$
  \State construct route-specific context and reward structures
\EndFor
\State initialize replay and $\pi$ over $\{\mathcal{X}_r(T_r)\}$
\State initialize the configured static or history-dependent availability process
\For{$t=1,\ldots,H$}
  \State prepare availability from the configured process and prior history
  \State $r_t\gets\pi.\Call{SelectRoute}{}$ \Comment{model-specific}
  \State $\mathbf{x}_t\gets\pi.\Call{SelectAction}{\mathcal{X}_{r_t}(T_{r_t})}$ \Comment{model-specific within-route action}
  \State $\mathcal{F}_t\gets\Call{Feedback}{r_t,\mathbf{x}_t,A_t}$
  \State $\pi.\Call{Update}{r_t,\mathbf{x}_t,\mathcal{F}_t}$ \Comment{model-specific channels and gates}
  \State update replay and route statistics
  \State record the selected route for subsequent reactive availability
  \If{allocator update is enabled and due at $t$}
    \State $\mathbf{T}\gets A.\Call{Update}{t,\text{route statistics}}$
    \State regenerate and propagate affected action/context/reward structures
  \EndIf
\EndFor
\State \Return efficiency, regret, robustness, and stability metrics
\end{algorithmic}
\end{algorithm}

\smallTitle{Execution scope and feedback contract}
Algorithm~\ref{alg:joint_loop} describes an interface, not a claim that every optional hook executed in every corpus. The accepted medium campaign uses fixed inter-route budgets; no adaptive allocator-update cadence is inferred from it. Exact adaptive-allocator cadence in the historical primary corpus remains provenance-pending. Environment-state transition intervals do not establish allocator cadence.

The feedback record $\mathcal{F}_t$ distinguishes continuous selected payoff, sampled route feedback, and whether a within-route update is applied. In the accepted medium implementation, GNeuralUCB, EXPNeuralUCB, and CPursuitNeuralUCB use availability-masked Bernoulli route feedback, while their NeuralUCB action learner receives the selected base payoff when the route is available. EXP additionally importance-weights its route update. iCPursuitNeuralUCB instead supplies continuous selected payoff to its predictive route update; its absent recorded Bernoulli channel is not a sequence of zero observations. These channel differences are part of the evaluated model implementations, not an assertion that every policy learns from the same scalar target. The Bernoulli success model in \cref{subsec:reward} and the continuous medium evaluation outcome must therefore not be conflated.

\begin{table*}[t]
\centering\scriptsize
\setlength{\tabcolsep}{2pt}
\caption{Source-verified policy semantics for the validated portfolio. Route-selection behavior is reported from each concrete class rather than inferred from its family name.}
\label{tab:policy_semantics}
\begin{tabular}{p{0.19\textwidth}p{0.22\textwidth}p{0.25\textwidth}p{0.27\textwidth}}
\toprule
\textbf{Policy} & \textbf{Route selection} & \textbf{Within-route action selection} & \textbf{Update/information mechanism}\\
\midrule
\texttt{CEpsilonGreedy} & $\epsilon$-greedy over route means & route-specific $\epsilon$-greedy & selected route/action reward update\\
\texttt{CEXP4} & EXP4 & route-specific $\epsilon$-greedy & normalized expert advice; route/action updates\\
\texttt{CPursuit} & UCB over route means & route-specific Pursuit & empirical-best Pursuit update\\
\texttt{CEpochGreedy} & round-robin & route-specific EpochGreedy & selected-action update\\
\texttt{CThompsonSampling} & round-robin & route-specific Thompson sampling & selected-action posterior update\\
\texttt{iCEpsilonGreedy} & exploratory/best observed route mean & informed $\epsilon$-greedy & predictive iCMAB update\\
\texttt{iCEXP4} & random & informed EXP4 & normalized advice; predictive update\\
\texttt{iCPursuit} & random & informed Pursuit & predictive Pursuit update\\
\texttt{iCEpochGreedy} & random & informed EpochGreedy & predictive iCMAB update\\
\texttt{iCThompsonSampling} & random & informed Thompson sampling & predictive iCMAB update\\
\texttt{EXPNeuralUCB} & EXP3 & NeuralUCB & adversarial route weights plus neural update\\
\texttt{EXPUCB} & EXP3 & linear UCB & adversarial route weights plus linear update\\
\texttt{GNeuralUCB} & simple UCB & NeuralUCB & route statistics plus neural update\\
\texttt{CPursuitNeuralUCB} & Pursuit/CMAB & NeuralUCB & Pursuit route probabilities plus neural update\\
\texttt{iCPursuitNeuralUCB} & predictive iCMAB/Pursuit & NeuralUCB & ARIMA-based prediction and anomaly-aware filtering\\
\texttt{Oracle} & exhaustive reward reference & maximizing feasible action & no learning; evaluation reference\\
\bottomrule
\end{tabular}
\end{table*}

\begin{table*}[t]
\centering\footnotesize
\setlength{\tabcolsep}{2pt}
\caption{Verified context representations. Untraced external-testbed encodings remain provenance-pending.}
\label{tab:context_representation}
\begin{tabular}{p{0.18\textwidth}p{0.38\textwidth}p{0.12\textwidth}p{0.25\textwidth}}
\toprule
\textbf{Scope} & \textbf{Representation} & \textbf{Dimension} & \textbf{Boundary}\\
\midrule
Primary matched evaluation & feasible per-link qubit split $\mathbf{x}\in\mathcal{X}_r(T_r)$ & $h_r$ (2 or 3) & route identity is structural; threat and history enter feedback/update\\
Medium anchor & nine qubits split over three links; 55 actions per route & 3 & ten route-specific sets, 550 route/action pairs\\
Paper8 external testbed & capacity, hops, min/mean fidelity, min/mean rate, minimum purification rounds, swap-success product & 8 & one within-path action vector per path\\
Other external testbeds & testbed-specific construction & provenance pending & no universal context encoding asserted\\
\bottomrule
\end{tabular}
\end{table*}

For the verified primary NeuralUCB implementation, a $d_r\!\to\!128\!\to\!1$ ReLU network scores each feasible action as
\[
s_t(\mathbf{x})=f_\theta(\mathbf{x})+\beta\sqrt{g_\theta(\mathbf{x})^\top\Sigma_t^{-1}g_\theta(\mathbf{x})}.
\]
The learner stores observations in replay, updates the gradient-based confidence state, and trains with Adam using mean-squared error plus squared parameter-drift regularization. The implementation begins training after $T>K$, performs two optimization steps per update, and samples replay minibatches of 64. These are verified mechanism details, not claims that current defaults equal every validated-run configuration.

For Pursuit, probabilities begin uniformly and move toward the empirical-best arm $i^\star$:
\[
p_{i^\star}\leftarrow p_{i^\star}+\alpha(1-p_{i^\star}),\qquad p_i\leftarrow p_i-\alpha p_i\quad(i\neq i^\star).
\]
Here $\alpha$ controls the movement rate. Reported-run values for $\alpha$, NeuralUCB confidence/regularization parameters, predictive settings, anomaly thresholds, and replay overrides remain provenance-pending unless tied to the validated corpus.

\begin{table}[t]
\centering\footnotesize
\caption{Settings explicitly recorded in all 75 accepted medium manifests. These are not substituted for unresolved historical run values.}
\label{tab:medium_recorded_settings}
\begin{tabular}{@{}p{0.36\linewidth}p{0.56\linewidth}@{}}
\toprule
Item & Recorded value\\
\midrule
Allocator/budget & Default/fixed; 90 total, nine per route\\
Horizon/blocks & 6,000 frames; three blocks\\
Replay & $T_b$; base 6,000; scale 2; capacity 12,000\\
Base seed & 12,345; component seeds derived by recorded SHA-256 rule\\
G / CP & neural mode; $\beta=1$\\
EXP & hybrid mode; $\beta=1$; eta factor 0.05; gamma factor 0.01\\
iCP & neural mode; $\beta=1$; eta factor 0.005; gamma factor 0.1; observation override unset\\
Oracle & base mode; privileged comparator\\
\bottomrule
\end{tabular}
\end{table}

The medium manifests pin the execution revision and imported-source hashes, and distinguish explicit keyword arguments from constructor defaults. \Cref{tab:medium_recorded_settings} reports explicit recorded settings only. The neural architecture/training description above is an implementation mechanism specification; unverified effective nested-learner settings are not promoted to reported-run hyperparameters. Historical Pursuit rates, neural training overrides, predictive orders/windows, and external encodings other than Paper8 remain separate provenance items. A class existing in the framework does not establish its inclusion in a corpus: the medium roster has four learned policies plus Oracle, whereas the historical portfolio has fifteen non-Oracle policies.

\subsection{Pinned Medium Availability Processes}
\label{app:medium_threats}
The accepted medium implementation records no added interruption for Baseline and independent path/frame interruption probability 0.0625 for Stochastic. Markov uses an independent binary state per path, initialized with attacked probability 0.25; each frame retains the state with probability 0.7 and otherwise toggles it. Thus 0.25 is an initialization parameter, not its long-run interruption fraction.

Adaptive uses the previous 100 selected routes. If $f_r$ is a route's frequency in that available history, its interruption probability is $\min(0.25+0.5f_r,0.9)$; with empty history it is 0.25. OnlineAdaptive starts from pending burst state, initiates an all-route burst with probability 0.3 for up to ten frames, and otherwise, from frame index five onward, targets the most recently selected route with probability 0.5. That targeted assignment can overwrite a pending-burst zero; routes left available are then independently interrupted with probability 0.25. This ordering is the pinned executed rule, not an idealized attack model. It explains why a configured base rate cannot be equated with measured global availability.

Static masks match within each medium block/threat; reactive masks depend on the individual policy's prior route history. The shared seeds therefore do not guarantee a common reactive counterfactual. These processes instantiate simulator dependence/reactivity, not measured physical attack distributions or a monotonic severity ladder. They do not establish the unresolved historical parameter values in \cref{tab:scenario_specs}.
